% paper/sections/appendix_ccd.tex
%
% 付録C: CCD 法の係数導出と実装（旧付録B+C を統合）
%   — §4/§8 対応

\section{CCD 法の係数導出と実装}
\label{app:ccd_supplement}

%% C.1 CCD 係数の Taylor 展開導出（内点＋境界）
\subsection{CCD 係数の Taylor 展開導出}
\label{app:ccd_coef_all}
\input{sections/appendix_ccd_coef_s1}    %% C.1.1 内点係数
\section{CCD 境界スキームの係数導出}
\label{app:ccd_bc_derivation}
% ═══════════════════════════════════════════════════════════════════════════════

本付録は §\ref{sec:ccd_bc} の左境界スキームの係数を Taylor 展開により導出する．

\subsection{Equation-I 境界スキームの導出}
\label{app:ccd_bc_derivation_I}

境界点 $i=0$ において，片側スキーム（式 \eqref{eq:bc_general}）の未知数
$(\alpha, \beta, c_0, c_1, c_2, c_3)$ を決定する．

\paragraph{Step 1：$\alpha, \beta$ の決定}

$\fp{1}, \fpp{1}$ を $i=0$ 周りで展開して LHS に代入すると，
$\fpp{0}$ の係数消去条件 $\alpha + \beta = 0$（すなわち $\beta = -\alpha$）が得られる．

\paragraph{Step 2：$\Ord{h^5}$ 精度から $\alpha$ の一意決定}

$\beta = -\alpha$ 代入後，$S_k \equiv c_1 + 2^k c_2 + 3^k c_3$ として
5つの整合条件（定数項から $f^{(4)}_0$ まで）を立てる：
$c_0{+}c_1{+}c_2{+}c_3=0$，$S_1=1{+}\alpha$，$S_2=0$，$S_3=-3\alpha$，$S_4=-8\alpha$．

差分 $S_2{-}S_1$，$S_3{-}S_2$，$S_4{-}S_3$ を順次取ると：
\[
  2c_2+6c_3 = -(1+\alpha),\quad
  4c_2+18c_3 = -3\alpha,\quad
  8c_2+54c_3 = -5\alpha
\]
第1式の2倍を第2式から引くと $c_3 = (2-\alpha)/6$．
第1式に戻して $c_2 = -3/2$（$\alpha$ に依らず一定）．
第3式に代入：
\[
  -12 + 9(2-\alpha) = -5\alpha
  \quad\Longrightarrow\quad
  6 = 4\alpha
  \quad\Longrightarrow\quad
  \boxed{\alpha = \tfrac{3}{2}}
\]
$\alpha \neq 3/2$ では $\Ord{h^4}$ にとどまる—— $\alpha = 3/2$ は唯一の $\Ord{h^5}$ 選択である．

\paragraph{Step 3：残余係数の計算}

$\alpha = 3/2$，$c_2 = -3/2$ を用いて：
$c_3 = 1/12$，$c_1 = 21/4$，$c_0 = -23/6$．\qed

\subsection{Equation-II 境界スキームの導出}
\label{app:ccd_bc_derivation_II}

境界点 $i=0$ において，式 \eqref{eq:bcII_general} の未知数
$(\gamma, \delta, d_0, d_1, d_2, d_3)$ を決定する．

\paragraph{Step 1：$\gamma, \delta$ の決定}

$\fp{1}, \fpp{1}$ を $i=0$ 周りで展開して LHS に代入すると：
\[
  \text{LHS} = \gamma h\,\fp{0} + (1+\gamma+\delta)\fpp{0} + \cdots
\]
$\fp{0}$ の係数をゼロにする：$\gamma = 0$．
$\fpp{0}$ の係数を 1 に正規化：$1+\gamma+\delta = 1$，すなわち $\delta = 0$．
したがって LHS $= \fpp{0}$（片側関数値スキーム）．

\paragraph{Step 2：RHS の係数決定}

$f_k$（$k=0,1,2,3$）を $f_0$ 周りで展開し，$f_0/h^2$，$\fp{0}/h$，$\fpp{0}$，$f^{(3)}_0/h$ の
各係数整合条件を解く：
\begin{itemize}
  \item $\fpp{0}$係数$-$$\fp{0}/h$係数: $d_2+3d_3=1$
  \item $f^{(3)}_0/h$係数$-$$\fpp{0}$係数: $2d_2+9d_3=-1$
  \item $3d_3 = -3 \Rightarrow d_3=-1$，$d_2=4$，$d_1=-5$，$d_0=2$
\end{itemize}
結果：式 \eqref{eq:bcII_left} の $\fpp{0} = (2f_0-5f_1+4f_2-f_3)/h^2$（$\Ord{h^2}$ 精度）．\qed

\subsection{Equation-II 境界スキームの4次精度昇格（\texorpdfstring{$\Ord{h^4}$}{O(h\textasciicircum 4)}）}
\label{app:ccd_bc_derivation_II_h4}

4点公式（$\Ord{h^2}$）では全体の CCD d2 収束が境界で制限されるため，6点公式
（$\Ord{h^4}$）に昇格する（$n_\mathrm{pts} \ge 6$ のとき有効）．

\paragraph{導出}

$f_k$（$k=0,\ldots,5$）の Taylor 展開を $f_0$ 周りで行い，係数
$d_0,\ldots,d_5$ を連立方程式
\[
  \sum_{k=0}^{5} d_k \, f_k = \fpp{0} + \Ord{h^p}, \qquad p \ge 4
\]
の条件 $k=0,1,2,3,4$ から決定する（$k=5$ は $\Ord{h^4}$ 打ち切り誤差の確認）：

\begin{align}
  k=0:\quad & \sum d_k = 0         \label{eq:bcII_h4_k0}\\
  k=1:\quad & \sum k\,d_k = 0      \label{eq:bcII_h4_k1}\\
  k=2:\quad & \sum k^2 d_k = 2     \label{eq:bcII_h4_k2}\\
  k=3:\quad & \sum k^3 d_k = 0     \label{eq:bcII_h4_k3}\\
  k=4:\quad & \sum k^4 d_k = 0     \label{eq:bcII_h4_k4}\\
  k=5:\quad & \sum k^5 d_k \ne 0   \quad (\text{$\Ord{h^4}$ 確認})
\end{align}

唯一解：
\begin{equation}
  \fpp{0}
  = \frac{45f_0 - 154f_1 + 214f_2 - 156f_3 + 61f_4 - 10f_5}{12h^2}
  + \Ord{h^4}.
  \label{eq:bcII_left_h4}
\end{equation}

\paragraph{数値検証}

係数ベクトル $\bm{d} = [45, -154, 214, -156, 61, -10]/12$ に対して，
$\sum_{k=0}^{5} d_k k^p$ が $p=0,\ldots,4$ ですべて $0$（$p=2$ は $2$）を与え，
$p=5$ で $-548/12 \ne 0$ となることを確認した（機械精度）．

\paragraph{実装}

\texttt{ccd/ccd\_solver.py} の \texttt{\_boundary\_coeffs\_left(h, n\_pts)} において，
$n_\mathrm{pts} \ge 6$ のとき式~\eqref{eq:bcII_left_h4} を，$n_\mathrm{pts} < 6$ のとき
旧来の $\Ord{h^2}$ 公式をフォールバックとして使用する．右境界では
$f_0 \to f_N,\, f_1 \to f_{N-1},\ldots$ の鏡像公式を適用する．

\paragraph{収束効果}

昇格により d2 の全域 $L^\infty$ 収束が境界寄与によって制限されず，
内点スキームの $\Ord{h^6}$ 精度に近い収束率が得られる．
数値実験（\texttt{test\_ccd.py} Test~2）では平均スロープ ${\approx}4.8$ を確認した．\qed

    %% C.1.2 境界スキーム係数

%% C.2 境界処理の代替実装（ゴーストセル法＋周期BC）
\subsection{境界処理の代替実装}
\label{app:ccd_boundary_alt}
\section{ゴーストセル法による CCD 境界処理（代替実装）}
\label{app:ccd_ghost}
% ═══════════════════════════════════════════════════════════════════════════════

本付録は §\ref{sec:ccd_bc} の境界スキーム法に対する代替として，\textbf{ゴーストセル（ghost cell）法}
を説明する．本稿の標準実装は境界スキーム法であり（§\ref{sec:ccd_bc}），
第\ref{sec:pressure}章（§\ref{sec:ccd_neumann_bc}）の Neumann 境界条件の組み込みも境界スキーム法を前提としている．
独自実装としてゴーストセル法を選択する場合は，PPE の Neumann 境界処理（§\ref{sec:ccd_neumann_bc}）を
ゴーストセル法に合わせて再設計すること．

計算領域の外側に仮想セル（$i = -1$，$i = N$）を1層追加し，
境界条件に応じた値を設定することで，内点スキーム（3点ステンシル）を境界近傍でもそのまま適用できる．
これにより行列アセンブリのコードが統一される利点がある．

\begin{tcolorbox}[algbox, title={ゴーストセルの設定手順（左境界 $i=0$ の例）}]
\textbf{Step 1：ゴーストセル値の設定}

境界条件に応じて $f_{-1}$，$f'_{-1}$，$f''_{-1}$ を以下のように設定する：

\begin{center}
\renewcommand{\arraystretch}{1.4}
\begin{tabular}{lll}
\hline
境界条件 & ゴーストセル値 & 物理的意味 \\
\hline
Dirichlet（ノースリップ壁）: $u|_{\mathrm{wall}}=0$ &
  $u_{-1} = -u_1$ &
  鏡像反転（$u_{0} \approx 0$） \\
Dirichlet（強制流入）: $u|_{\mathrm{wall}}=U_w$ &
  $u_{-1} = 2U_w - u_1$ &
  $u_{0} = (u_{-1}+u_1)/2 = U_w$ \\
Neumann（流出・対称）: $\partial u/\partial x|_{\mathrm{wall}}=0$ &
  $u_{-1} = u_1$ &
  勾配ゼロ（偶関数的反射） \\
\hline
\end{tabular}
\end{center}

微分値のゴーストセル設定（ノースリップ壁の例）：
$f'$ は偶関数的に反射（$f'_{-1} = f'_1$），$f''$ は奇関数的（$f''_{-1} = -f''_1$）．

\textbf{Step 2：境界点での内点方程式の適用}

\begin{equation}
  \alpha_1 f'_{-1} + f'_0 + \alpha_1 f'_1
  = \frac{a_1}{h}(f_1 - f_{-1}) + b_1 h(f''_1 - f''_{-1})
  \label{eq:ccd_ghost_eq}
\end{equation}
ここで $f_{-1}$，$f'_{-1}$，$f''_{-1}$ はStep 1 で設定したゴーストセル値．
全格子点で同一の係数行列 $\mathbf{A}_L, \mathbf{A}_R, \mathbf{B}$ を使えるため，行列アセンブリが単純化される．

\textbf{注意：} 圧力の全 Neumann 境界問題（零固有モード）については第\ref{sec:pressure}章の対処法を参照すること．
\end{tcolorbox}

\begin{tcolorbox}[algbox, title={PPE Neumann 境界条件のゴーストセル法による再設計（$\partial p/\partial n = 0$ の場合）}]
境界スキーム法（§\ref{sec:ccd_bc}，§\ref{sec:ccd_neumann_bc}）では，
第1行を $p'_0 = 0$ に直接置き換えることで Neumann 条件を強制した．
ゴーストセル法を選択した場合，この条件は\textbf{ゴーストセル値の偶関数的反射}として実装する：

\medskip
\textbf{Step 1：圧力ゴーストセルの設定（左境界 $i=0$ の例）}
\[
  p_{-1} = p_1, \qquad p'_{-1} = p'_1, \qquad p''_{-1} = p''_1
\]
（偶関数的反射：$\partial p/\partial x|_{x=0} = 0$ を意味する）

\textbf{Step 2：境界点での内点方程式の適用}

ゴーストセル値を代入した Equation-I（式~\eqref{eq:ccd_ghost_eq}）：
\[
  \alpha_1 p'_{-1} + p'_0 + \alpha_1 p'_1
  = \frac{a_1}{h}(p_1 - p_{-1}) + b_1 h(p''_1 - p''_{-1})
\]
$p_{-1} = p_1$ および $p''_{-1} = p''_1$ を代入すると，右辺が消え：
\[
  (1 + 2\alpha_1) p'_0 = 0 \quad\Longrightarrow\quad p'_0 = 0
\]
Neumann 条件が自動的に成立する（境界スキームの明示的な修正は不要）．

\textbf{Step 3：CCD 演算子値の計算}

$p'_0 = 0$ が確定したうえで，変密度 PPE 演算子値
$(\mathcal{L}_{\mathrm{CCD}}^{\rho,x}\bm{p})_0 = p''_0/\rho_0$
（§\ref{sec:ccd_neumann_bc}，式の導出と同様；密度勾配項は $p'_0=0$ により消滅）を計算する．

\medskip
\textbf{零固有モードへの対処：} ゴーストセル法を使う場合も，全周 Neumann 境界での零モード問題（§\ref{sec:pinpoint}）は同様に発生する．式~\eqref{eq:pin_overwrite} の強制上書きによるピン条件が有効であり，ゴーストセル法特有の追加対処は不要である．
\end{tcolorbox}

% ═══════════════════════════════════════════════════════════════════════════════
    %% C.2.1 ゴーストセル法
\input{sections/appendix_ccd_impl_s4}    %% C.2.2 周期境界条件

%% C.3 精度解析と楕円型ソルバー（混合微分精度＋PPEソルバー構造＋Kronecker積）
\subsection{精度解析と楕円型ソルバー}
\label{app:ccd_analysis}
\subsection{2 段階 CCD 混合微分の精度解析：\texorpdfstring{$C^8$}{C\textasciicircum 8} 条件と劣化機構}
\label{app:ccd_mixed_precision}
% ═══════════════════════════════════════════════════════════════════════════════

§\ref{sec:ccd_mixed} の2段階 CCD 混合微分 $\phi_{xy}$ の精度を詳細に解析する．

Step 1 で得た $g = \partial f/\partial x$ の誤差は $E_g = O(h^6)$．
Step 2 で $g$ を $y$ 方向に微分するとき，誤差の上界は：
\[
  E_{\phi_{xy}} \leq \underbrace{O(h^6)}_{\text{$y$ 方向 CCD の打ち切り誤差}}
               + \underbrace{\left\|\pd{E_g}{y}\right\|}_{\text{$E_g$ の $y$ 方向変動}}
\]

\textbf{一般論では $\partial E_g/\partial y \sim O(h^6/h) = O(h^5)$ となる可能性がある．}
差分演算子は入力の誤差を $h^{-1}$ 倍程度増幅するため，$O(h^6)$ の誤差場を数値微分すると $O(h^5)$ に劣化するのが一般論である．

\textbf{ただし $f \in C^8$ であれば $O(h^6)$ が維持される：}
CCD の打ち切り誤差 $E_g(x,y) \propto h^6 f^{(7)}$ は $f$ の高次微分から生じる滑らかな場であり，
その $y$ 方向変動 $\partial E_g/\partial y \propto h^6 f^{(8)}$ は $f \in C^8$ のもとで $O(h^6)$ オーダーにとどまる．

\textbf{精度まとめ：}
\begin{itemize}
  \item $f \in C^8$ の場合：$E_{\phi_{xy}} = O(h^6)$ が維持される
  \item $f$ に不連続性や急峻な変化がある場合：精度が $O(h^5)$ 以下に劣化する可能性あり
\end{itemize}

\textbf{気液二相流の界面近傍への影響：}
気液二相流では，界面近傍で CLS 変数 $\psi$ が $O(\varepsilon)$ の厚さで急峻に変化するため，
$\psi$ の高次微分（$\psi^{(8)}$ 等）は $O(1/\varepsilon^7)$ のオーダーとなり，
実質的に $C^8$ 条件が界面近傍の格子点（界面から $3\varepsilon$ 以内）で成立しない．
したがって，曲率計算の混合微分 $\phi_{xy}$ の精度は界面近傍で $O(h^5)$ 以下に劣化しうる．
ただし，この劣化は界面近傍の $O(\varepsilon/h) \sim O(1)$ 個の格子点のみに影響し，
大域的な $L^2$ ノルム誤差への寄与は $O(h^{5+d/2})$（$d$：次元）程度に収まる
（比較：CSF モデル誤差 $O(\varepsilon^2) \approx O(h^2)$ が既に界面近傍精度を律速している）．
実務上は，格子収束テスト（第\ref{sec:verification}章）で曲率 $\kappa$ の収束次数を直接確認することが
この劣化の存在と程度を評価する最も確実な方法である．

実証的確認として，$f(x,y) = \sin(\pi x)\cos(\pi y)$（解析解 $f_{xy} = -\pi^2\cos(\pi x)\sin(\pi y)$）で格子収束テストを行い，$N$ を倍にするたびに誤差が $1/64$ 倍（6次収束）になるかを確認すること（第\ref{sec:verification}章参照）．

% ═══════════════════════════════════════════════════════════════════════════════
    %% C.3.1 混合微分精度解析
\section{CCD 法の楕円型ソルバーとしての双対的役割}
\label{app:ccd_elliptic}
% ═══════════════════════════════════════════════════════════════════════════════

本付録は §\ref{sec:ccd_mixed} までで説明した「微分演算子としての CCD」に加えて，
CCD が\textbf{楕円型方程式の陰的ソルバー}として機能する数学的構造を説明する．
これは第\ref{sec:pressure}章の CCD-Poisson ソルバー（§\ref{sec:ppe_pseudotime}）の理論的基盤である．

\subsection{CCD 演算子行列の定義}

既知の関数値 $\bm{f} = [f_0, \dots, f_{N-1}]^\top$ に対してブロック Thomas 法で
$\bm{v} = [f'_0, f''_0, \dots, f'_{N-1}, f''_{N-1}]^\top$ を求める系：
\begin{equation}
  \mathcal{M}_{\mathrm{CCD}}\,\bm{v} = \mathcal{R}(\bm{f})
  \label{eq:ccd_derivative_system_app}
\end{equation}
から，第1・第2微分成分を抽出する射影行列 $\mathbf{E}_1, \mathbf{E}_2 \in \mathbb{R}^{N \times 2N}$ を用いると，
CCD 離散微分演算子は：
\begin{align}
  D_x^{(1)}\bm{p} &\equiv \mathbf{E}_1\,\mathcal{M}_{\mathrm{CCD}}^{-1}\,\mathcal{R}(\bm{p}) \\
  D_x^{(2)}\bm{p} &\equiv \mathbf{E}_2\,\mathcal{M}_{\mathrm{CCD}}^{-1}\,\mathcal{R}(\bm{p})
\end{align}
となる．1次元等密度ラプラシアンを離散化した演算子行列：
\begin{equation}
  \mathbf{L}_{\mathrm{CCD}}^{(x)} \equiv \mathbf{E}_2\,\mathcal{M}_{\mathrm{CCD}}^{-1}\,\mathcal{R}
  \in \mathbb{R}^{N \times N}
  \label{eq:L_CCD_def_app}
\end{equation}
は $N \times N$ の行列であり，$\mathbf{L}_{\mathrm{CCD}}^{(x)}\bm{p} \approx [p''_0, \dots, p''_{N-1}]^\top$（$\Ord{h^6}$ 精度）を返す．
$\mathbf{L}_{\mathrm{CCD}}^{(x)}$ を陽に組み立てる必要はなく，各反復でブロック Thomas 法を適用する matrix-free 実装が可能（$O(N)$ コスト）．

\subsection{変密度場における CCD-Poisson 演算子}

変密度 PPE の演算子 $\nabla\cdot(1/\rho\,\nabla p)$ を1次元で展開すると：
\[
  \mathcal{L}^{\rho}p = \frac{1}{\rho}\,p'' - \frac{\rho_x}{\rho^2}\,p'
\]
格子点ごとの密度行列 $\boldsymbol{\Lambda}_\rho = \mathrm{diag}(\rho_0, \dots, \rho_{N-1})$ を用いると：
\begin{equation}
  \mathbf{L}_{\mathrm{CCD}}^{\rho,x}\bm{p}
  = \boldsymbol{\Lambda}_\rho^{-1}\,\mathbf{L}_{\mathrm{CCD}}^{(x)}\bm{p}
  - \boldsymbol{\Lambda}_\rho^{-2}\,\mathrm{diag}(D_x^{(1)}\bm{\rho})\,D_x^{(1)}\bm{p}
\end{equation}
2次元ではこれを $x$・$y$ 方向に加算する．

等密度場では $\mathbf{L}_{\mathrm{CCD}}^{(x)}$ は適切な境界条件のもとで\textbf{負半定値}（全固有値 $\lambda_k \le 0$），ラプラシアンの本質的性質と一致する．
変密度場では厳密な対称性は失われるが，仮想時間反復の収束性はスペクトル半径に支配されるため実用上は十分である．

\subsection{PPE の数値解法：反復法（LGMRES）優先・直接 LU フォールバック}
\label{app:ccd_lu_direct}

PPE $\mathbf{L}_{\mathrm{CCD}}^{\rho}\bm{p} = \bm{q}$
（クロネッカー積表現：式~\eqref{eq:L_CCD_2d_kron}，付録~\ref{app:ccd_kronecker}）の求解には
\textbf{二段構えの戦略}を採用する．

\subsubsection*{Primary：LGMRES 反復法}

LGMRES（拡張 Krylov 部分空間法；\texttt{scipy.sparse.linalg.lgmres}）を第一選択とする．
収束した場合のメモリコストは $O(n k)$（$k$：内部再スタート数）であり，
行列要素のみを格納する疎行列表現と合わせて全体で $O(n)$ に抑えられる．
前タイムステップの圧力 $\bm{p}^n$（または増分法では $\delta\bm{p}^n$）を初期推定値として
warm start することで，収束に要するイテレーション数を大幅に削減できる．

LGMRES は通常の GMRES より拡張 Krylov 部分空間を使用するため，
CCD 境界スキームが誘発する高非対称行列（$N=16$ 時点で最大非対称度 $\approx 900$）でも
より頑健に収束することが期待される．
収束判定は許容誤差 $\varepsilon_{\mathrm{tol}}$（設定値 \texttt{pseudo\_tol}）と
最大反復回数 \texttt{pseudo\_maxiter} による．

\subsubsection*{Fallback：直接スパース LU 法}

LGMRES が \texttt{pseudo\_maxiter} ステップ以内に収束しない場合，
直接スパース LU 法（SuperLU: \texttt{scipy.sparse.linalg.spsolve}）に自動的に切り替える．
直接法はメモリコスト $O(n^{1.5})$（2次元問題での LU 充填の経験則）と引き換えに，
非特異な系を必ず正確に解くことが保証される．
この安全網により，反復法の収束失敗が他モジュールの開発を停滞させることを防ぐ．

\subsubsection*{零空間のピン条件}

等密度場で全周 Neumann 境界を課した PPE は定数圧力モードという零空間を持つ（§\ref{app:checkerboard_mode}）．
変密度場でも同様に零空間が生じうる（詳細は付録~\ref{app:ccd_kronecker}）．
零空間を除去するため，\textbf{中央ノード $(N/2, N/2)$}（整数除算）の行を恒等行で上書きして
その圧力を $0$ に固定する（ゲージ条件；式~\eqref{eq:pin_overwrite}）：
修正系 $\hat{\mathbf{L}}\bm{p} = \hat{\bm{q}}$ を LGMRES（または LU）で解く．
中央ノードは正方形領域の4つの対称操作（$x$・$y$ 軸反転，対角置換）に対して不変であり，
ピン固定がシミュレーションの物理的対称性を破らない（コーナー固定では対称性が崩れる）．

\subsubsection*{仮想時間フレームワークとの関係}

付録~\ref{sec:dtau_derive}（収束率解析）は，スウィープ型仮想時間反復の最適擬似時間刻み
$\Delta\tau_{\mathrm{opt}}$ を導出しており，数学的証明として引き続き有効である．
本付録の実装はスウィープ分割を行わず，組み立てた大域行列に対して直接 Krylov 法を適用するため，
スプリッティング誤差（付録~\ref{sec:dtau_derive} の $\mathcal{O}(\Delta\tau^2)$ 項）は生じない．

% ═══════════════════════════════════════════════════════════════════════════════
\subsection{2次元 CCD-Poisson 演算子のクロネッカー積による大域行列表現}
\label{app:ccd_kronecker}
% ═══════════════════════════════════════════════════════════════════════════════

本節では，§\ref{sec:ccd_2d_full} の点ごと演算子式（式~\eqref{eq:L_CCD_2d_full}）を，
付録~\ref{app:ccd_lu_direct} の直接 LU 法に渡す $n \times n$ スパース行列として
\textbf{代数的に厳密}に組み立てる手順を示す．

\subsubsection{記号と格子設定}

$x$ 方向に $N_x$ 点，$y$ 方向に $N_y$ 点の2次元等間隔ノードセンタード格子を考える（各端点を含む）：
\[
  N_x = N^{(x)}+1,\quad N_y = N^{(y)}+1,\quad n = N_x N_y
\]
2次元配列 $p[i,j]$（$0 \le i < N_x$，$0 \le j < N_y$）を \textbf{C順序（行優先）}でフラット化する：
\begin{equation}
  k(i,j) = i N_y + j \in \{0, 1, \dots, n-1\}
  \label{eq:c_order_index}
\end{equation}
フラットベクトル $\bm{p} \in \mathbb{R}^n$ を $\bm{p}_k = p[i,j]$ と定義する．

\subsubsection{1次元 CCD 行列の構成}

各軸の1次元 CCD 微分行列を以下のように定義する：
\begin{align}
  \mathbf{D}_1^{(x)} \in \mathbb{R}^{N_x \times N_x}: &\quad
    [\mathbf{D}_1^{(x)}]_{i\kappa} = \bigl[D_x^{(1)}\mathbf{e}_\kappa\bigr]_i
    \label{eq:D1x_def} \\
  \mathbf{D}_2^{(x)} \in \mathbb{R}^{N_x \times N_x}: &\quad
    [\mathbf{D}_2^{(x)}]_{i\kappa} = \bigl[D_x^{(2)}\mathbf{e}_\kappa\bigr]_i
    \label{eq:D2x_def}
\end{align}
ここで $\mathbf{e}_\kappa \in \mathbb{R}^{N_x}$ は第 $\kappa$ 標準基底ベクトル，$D_x^{(\ell)}$ は
§\ref{app:ccd_elliptic} のブロック Thomas 解法による CCD 微分演算子である．
$y$ 方向も同様に $\mathbf{D}_1^{(y)}, \mathbf{D}_2^{(y)} \in \mathbb{R}^{N_y \times N_y}$ を定義する．

\noindent\textbf{実装：} $\mathbf{D}_\ell^{(x)}$ は $N_x \times N_x$ 単位行列を入力として CCD を一括適用することで
$O(N_x^2)$ で構成できる（全列を一括差分）：
\[
  [\mathbf{D}_1^{(x)},\, \mathbf{D}_2^{(x)}] = \mathrm{CCD.differentiate}(\mathbf{I}_{N_x},\;\mathrm{axis}=0)
\]

\subsubsection{クロネッカー積による2次元微分行列}

C順序フラットベクトル $\bm{p}$ への2次元偏微分はクロネッカー積で表現される．
以下に $x$ 方向の場合の厳密な証明を示す．

\textbf{命題：} C順序のもとで，$(D_x^{(2)} p)_{ij}$ のフラット表現は $(\mathbf{D}_2^{(x)} \otimes \mathbf{I}_{N_y})\bm{p}$ に等しい．

\textbf{証明：}
クロネッカー積 $\mathbf{D}_2^{(x)} \otimes \mathbf{I}_{N_y}$ の $k(i,j)$ 行への作用は：
\[
  \bigl[(\mathbf{D}_2^{(x)} \otimes \mathbf{I}_{N_y})\bm{p}\bigr]_{i N_y + j}
  = \sum_{\kappa=0}^{N_x-1}\sum_{l=0}^{N_y-1}
    [\mathbf{D}_2^{(x)}]_{i\kappa}\,[\mathbf{I}_{N_y}]_{jl}\,\bm{p}_{\kappa N_y + l}
  = \sum_{\kappa=0}^{N_x-1} [\mathbf{D}_2^{(x)}]_{i\kappa}\,p[\kappa, j]
  = (D_x^{(2)} p)_{ij}
  \quad\square
\]
$y$ 方向も同様に，$(\mathbf{I}_{N_x} \otimes \mathbf{D}_2^{(y)})\bm{p}$ の $k(i,j)$ 成分は：
\[
  \bigl[(\mathbf{I}_{N_x} \otimes \mathbf{D}_2^{(y)})\bm{p}\bigr]_{i N_y + j}
  = \sum_{l=0}^{N_y-1} [\mathbf{D}_2^{(y)}]_{jl}\,p[i, l]
  = (D_y^{(2)} p)_{ij}
  \quad\square
\]
4つの2次元微分行列を以下のように定める（サイズはすべて $n \times n$）：
\begin{align}
  \mathbf{D}_{2x}^{\mathrm{2D}} &= \mathbf{D}_2^{(x)} \otimes \mathbf{I}_{N_y}
  \label{eq:D2x_2d} \\
  \mathbf{D}_{2y}^{\mathrm{2D}} &= \mathbf{I}_{N_x} \otimes \mathbf{D}_2^{(y)}
  \label{eq:D2y_2d} \\
  \mathbf{D}_{1x}^{\mathrm{2D}} &= \mathbf{D}_1^{(x)} \otimes \mathbf{I}_{N_y}
  \label{eq:D1x_2d} \\
  \mathbf{D}_{1y}^{\mathrm{2D}} &= \mathbf{I}_{N_x} \otimes \mathbf{D}_1^{(y)}
  \label{eq:D1y_2d}
\end{align}

\subsubsection{変密度 2次元 CCD-Poisson 演算子の行列表現}

密度場 $\rho[i,j]$ および各軸の密度勾配 $(\partial_x\rho)[i,j]$，$(\partial_y\rho)[i,j]$ を
同じ C順序でフラット化した対角行列：
\begin{align}
  \boldsymbol{\Lambda}_{\rho^{-1}}       &= \mathrm{diag}(1/\bm{\rho})
  \in \mathbb{R}^{n \times n} \\
  \boldsymbol{\Lambda}_{\partial\rho_x/\rho^2} &= \mathrm{diag}(\partial_x\bm{\rho}/\bm{\rho}^2)
  \in \mathbb{R}^{n \times n} \\
  \boldsymbol{\Lambda}_{\partial\rho_y/\rho^2} &= \mathrm{diag}(\partial_y\bm{\rho}/\bm{\rho}^2)
  \in \mathbb{R}^{n \times n}
\end{align}
を用いると，変密度 2次元 CCD-Poisson 演算子の大域行列表現は：
\begin{equation}
  \boxed{
  \mathbf{L}_{\mathrm{CCD}}^{\rho}
  = \boldsymbol{\Lambda}_{\rho^{-1}}\!
    \bigl(\mathbf{D}_{2x}^{\mathrm{2D}} + \mathbf{D}_{2y}^{\mathrm{2D}}\bigr)
  - \boldsymbol{\Lambda}_{\partial\rho_x/\rho^2}\,\mathbf{D}_{1x}^{\mathrm{2D}}
  - \boldsymbol{\Lambda}_{\partial\rho_y/\rho^2}\,\mathbf{D}_{1y}^{\mathrm{2D}}
  \;\in \mathbb{R}^{n \times n}
  }
  \label{eq:L_CCD_2d_kron}
\end{equation}
この式は §\ref{sec:ccd_2d_full} の点ごと式（式~\eqref{eq:L_CCD_2d_full}）と代数的に等価であり，
C順序インデックス写像~\eqref{eq:c_order_index} のもとで厳密に成立する．

\textbf{物理的解釈：}
$\mathbf{D}_{2x}^{\mathrm{2D}} + \mathbf{D}_{2y}^{\mathrm{2D}}$ は等密度等方ラプラシアン $\nabla^2$
の $O(h^6)$ 離散化であり，$\boldsymbol{\Lambda}_{\rho^{-1}}$ が変密度係数 $1/\rho$ を乗じ，
クロス微分項 $\boldsymbol{\Lambda}_{\partial\rho/\rho^2}\mathbf{D}_1^{\mathrm{2D}}$ が積の法則
$\nabla\cdot(1/\rho\,\nabla p) = (1/\rho)\nabla^2 p - (\nabla\rho/\rho^2)\cdot\nabla p$
（式~\eqref{eq:varrho_product_rule}）の第2項を担う．

\subsubsection{零空間の次元とピン条件の必要性}

等密度・全周 Neumann 境界では $\mathbf{L}_{\mathrm{CCD}}^{\rho}$ の零空間は定数モード1次元のみとなる
（連続問題に対応）．しかし Eq-II-bc（§\ref{sec:ccd_bc}）の $O(h^2)$ 一側公式は
ブロック Thomas 系に非対称成分を導入し，数値的に零空間次元が増大することがある
（$N=4$ 時点で零空間次元 $= 8$，条件数 $\approx 1.9 \times 10^{17}$；付録~\ref{app:ccd_lu_direct} も参照）．

したがってピン条件（式~\eqref{eq:pin_overwrite}）は\textbf{必須}である．
根本的解決策は Eq-II-bc を $O(h^4)$ 以上の公式に昇格させることであり，
これにより零空間次元が理論通り1次元に収束する（今後の課題；§\ref{sec:conclusion}）．

\subsubsection{計算コスト}

\begin{center}
\renewcommand{\arraystretch}{1.2}
\begin{tabular}{llp{5.5cm}}
\hline
処理 & コスト & 備考 \\
\hline
1D CCD 行列構成 & $O(N_x^2 + N_y^2)$ & 単位行列一括差分（タイムステップ毎ではなく初期化時1回） \\
クロネッカー積組み立て & $O(N_x N_y^2 + N_x^2 N_y)$ & scipy.sparse.kron \\
密度重みの対角行列更新 & $O(n)$ & 各タイムステップで $\rho$ 更新時に再計算 \\
スパース LU 因数分解 & $O(n^{1.5})$（2D 問題の経験則） & 各タイムステップ1回 \\
前進後退代入（求解） & $O(n^{1.0})$ & LU 因数後 \\
\hline
\end{tabular}
\end{center}

典型的な計算では $n \lesssim 10^4$（$100 \times 100$ 格子）のため，
1タイムステップあたりの PPE 求解コストは $O(n^{1.5}) \approx 10^6$ 演算程度であり実用的である．

% ═══════════════════════════════════════════════════════════════════════════════
    %% C.3.2 楕円型ソルバー・Kronecker積・Neumann BC
